\section*{Ordbank --- Sentrale fagbegrep og teorier}
\addcontentsline{toc}{section}{Ordbank --- Sentrale fagbegrep og teorier}
\label{sec:ordbank}

Felles oppslagsverk for alle forelesningene og prosjektforsvaret.
Kortdefinisjoner du skal kunne aktivt --- ikke bare gjenkjenne.

\subsection*{A. Hovedteoriene oppgaven bygger på}

\subsubsection*{DVC Framework --- Digital Value Creation}
\textbf{Kilde:} Duus \& Cooray (2020), \emph{The European Business Review}.\\
Rammeverk for å diagnostisere \emph{hvor} og \emph{hvordan} digital kundeverdi
skapes. Tre indre lag (experiences, relationships, evolution) med
\term{relevance} som integrerende kjerne, og tre ytre forutsetninger (digital
infrastruktur, kompetanser, outputs).\\
\keypoint{Bruk:} Diagnostisere appens kundeverdi. Svakhet: kan overvurdere
aktivitet som verdi --- må kobles til DBM.

\subsubsection*{DBM --- Digital Business Model}
\textbf{Kilde:} Weill \& Woerner (2018), \emph{What's Your Digital Business
Model?}\\
Plasserer virksomheter langs to akser --- \emph{customer knowledge} og
\emph{business design} --- og gir fire arketyper (omnichannel, supplier,
ecosystem driver, modular producer).\\
\keypoint{Bruk:} Tester om digital verdi passer forretningsmodellen og om den
kan omsettes til \emph{value capture}.

\subsection*{B. DVC-komponentene}

\begin{tabular}{p{3.6cm}p{12cm}}
\hline
\textbf{Begrep} & \textbf{Definisjon} \\
\hline
Relevance & Kjerne i DVC. Om de digitale løsningene oppleves meningsfulle for
brukeren. Hele rammeverket roterer om dette. \\
\hline
Experiences & Hvordan digitale løsninger møter kundens behov, emotivt og
utilitaristisk. \\
\hline
Set experiences & Planlagte digitale opplevelser (handleliste, kvitteringer,
oppskrifter). \\
\hline
Situational experiences & Uventede situasjoner --- tester organisasjonens
reaksjonsevne (f.eks. Scan \& Pay-feil). \\
\hline
Relationships & Tillitsbaserte, personaliserte relasjoner via digitale kanaler;
kunden får økende kontroll. \\
\hline
Evolution & Hvordan eksisterende løsninger videreutvikles til nye tjenester,
modeller, partnerskap. \\
\hline
Digital infrastruktur & Back-end-kapabiliteter (API-lag, plattform,
lojalitetsmotor) som muliggjør digital verdi. \\
\hline
Digitale kompetanser & Organisatorisk evne til å bygge og bruke digitale
løsninger. \\
\hline
Digitale outputs & Målbare digitale resultater (MAU, klikk, innløsninger).
\emph{Mellomliggende} mål --- ikke utfall. \\
\hline
Outputs $\neq$ outcomes & Kritisk poeng: aktivitet er ikke det samme som
verdiskaping. Brukes for å forsvare KPI-hierarkiet (Anbefaling 1). \\
\hline
\end{tabular}

\subsection*{C. DBM-komponentene}

\subsubsection*{De fire arketypene}

\begin{tabular}{p{3.6cm}p{8.5cm}p{3.5cm}}
\hline
\textbf{Arketype} & \textbf{Kort definisjon} & \textbf{Coop-link} \\
\hline
Omnichannel\newline business & Fysisk + digital, høy kundekunnskap,
kontrollert verdikjede. & Coops hovedposisjon. \\
\hline
Supplier & Leverer til andres verdikjeder, svak kunderelasjon. & Risiko hvis
Coop kutter digitalt. \\
\hline
Ecosystem driver & Hub i bredt økosystem, høy kundekunnskap. & Ikke Coop
(Coop.dk lukket). \\
\hline
Modular producer & Plug-in-kapabilitet inn i andres økosystemer. & Lobyco
isolert sett. \\
\hline
\end{tabular}

\subsubsection*{Nøkkelbegreper}

\begin{tabular}{p{3.8cm}p{12cm}}
\hline
\textbf{Begrep} & \textbf{Definisjon} \\
\hline
Customer knowledge & Hvor godt virksomheten kjenner kundene --- fra delvis til
fullstendig. \\
\hline
Business design & Hvor kontrollert verdikjeden er --- fra integrert til
økosystem. \\
\hline
Value creation & Hvem som får verdi og hvordan (kunder, butikker, leverandører,
Coop). \\
\hline
Value capture & Hvordan virksomheten fanger økonomisk verdi (margin, retail
media, retensjon). Oppgavens kritiske gap. \\
\hline
Store-first\newline omnichannel & Coops posisjon: appen er et digitalt lag
\emph{rundt} det fysiske butikkbesøket, ikke en selvstendig kanal. \\
\hline
\end{tabular}

\subsection*{D. Andre pensumteorier}

\begin{tabular}{p{3.4cm}p{3.5cm}p{8.5cm}}
\hline
\textbf{Teori} & \textbf{Kilde} & \textbf{Kjerne} \\
\hline
Strategi som\newline posisjonering & Porter (1996) & Strategi = å utføre andre
aktiviteter enn konkurrentene, eller lignende aktiviteter på en annen måte.
Være \emph{annerledes}, ikke bare \emph{bedre}. \\
\hline
Porter's Five\newline Forces & Porter & Verktøy for ekstern bransjeanalyse:
rivalisering, kjøpermakt, leverandørmakt, substitutter, nyetablerere. \\
\hline
Verdikjede & Porter & Aktivitetene som skaper og leverer verdi til kunden. \\
\hline
Business Model\newline Canvas (BMC) & Osterwalder \& Pigneur (2010) &
Analytisk verktøy for å kartlegge forretningsmodellen i 9 byggesteiner. Ikke en
strategi i seg selv. \\
\hline
VRIO /\newline sustained advantage & Lektion 2 & Ressurs gir varig fordel hvis
den er Valuable, Rare, Imitable (vanskelig), Organized. \\
\hline
DT vs.\newline IT-Enabled\newline Transformation & Wessel et al. (2021) & DT:
digital teknologi redefinerer value proposition og identitet. ITE: teknologi
støtter eksisterende value proposition. \textbf{Coop = ITE.} \\
\hline
Digital\newline innovasjon & Yoo et al. (2010); Nambisan et al. (2017); Hukal
\& Henfridsson (2017) & Samskapelse av nye tilbud gjennom
\term{rekombinasjon} av digitale og fysiske komponenter. \\
\hline
Diffusion of\newline Innovations & Rogers (1962/1983) & Innovasjon =
ide/praksis/objekt som oppfattes som ny av en adopsjonsenhet. \\
\hline
Disrupsjon & Christensen & Disruptive aktører starter i bunnen/forsømt segment
med et enklere, billigere tilbud og forbedrer seg oppover. \\
\hline
Sustaining\newline innovation & Christensen & Forbedrer eksisterende produkter
for eksisterende kunder. \textbf{Coop App = sustaining}, ikke disruptiv. \\
\hline
Coop IT-legacy & Hedman \& Bjørn-Andersen (2016) & Coops historiske
teknologiske gjeld --- relevant for implementeringskapasitet. \\
\hline
3C-rammeverket & Tawse \& Tabesh (2021) & \emph{Competency, Commitment,
Coordination} --- tre forutsetninger for vellykket strategiimplementering. \\
\hline
Digitale\newline økosystemer & Lektion 6 & Nettverk av gjensidig avhengige
aktører som samskaper verdi. Plattform = hub. Drives av nettverkseffekter. \\
\hline
\end{tabular}

\subsection*{E. Egenskaper ved digital teknologi (Lektion 5)}

\begin{tabular}{p{4cm}p{12cm}}
\hline
\textbf{Egenskap} & \textbf{Betydning} \\
\hline
Datahomogenitet & Digital informasjon er universell ressurs --- kan deles og
kombineres på tvers. \\
\hline
Redigerbarhet & Digitalt innhold kan endres etter at det er produsert. \\
\hline
Selvrefererende\newline egenskaper & Digitale systemer kan referere til og
endre seg selv (rekursivitet). \\
\hline
Rekombinasjon & Digital innovasjon skjer ved å kombinere eksisterende digitale
+ fysiske komponenter på nye måter. \\
\hline
\end{tabular}

\subsection*{F. Begreper fra oppgaven og anbefalingene}

\begin{tabular}{p{3.8cm}p{12cm}}
\hline
\textbf{Begrep} & \textbf{Definisjon} \\
\hline
Coop App & Coops mobilapplikasjon --- digitalt lag rundt fysisk butikk. \\
\hline
Lobyco & Coops heleide datterselskap som drifter app-/lojalitetsplattformen og
selger den eksternt. \\
\hline
Retail media & Annonseringsinntekt fra leverandører som vil eksponeres mot
Coops kundebase i appen/butikken. \\
\hline
Scan \& Pay & Selvbetjent skanne-og-betal-funksjon i appen --- kjerne i set
experience i butikk. \\
\hline
Gamification & Spillmekanikker (poeng, utfordringer) brukt for å engasjere
kunder. Levert av Playable. \\
\hline
KPI-hierarki & Anbefaling 1: skille mellom drivere (MAU, klikk) og
forretningsutfall (kurv, retensjon, EBIT-bidrag). \\
\hline
Hybrid governance & Styringsmodell der konsernet eier plattform/KPI-er, mens
kjedene tilpasser innholdet. Brukes i Anbefaling 2. \\
\hline
Pilot-først & Implementeringsprinsipp: test små, isolerte tiltak før
utrulling, koblet til Lektion 7. \\
\hline
Customer Value\newline Proposition (CVP) & Hva kunden faktisk får av verdi
--- sentral i Lektion 3. \\
\hline
MAU & Monthly Active Users --- digital output, ikke outcome. \\
\hline
Household reach & Andel av husholdninger appen når --- leverandørclaim fra
Lobyco. \\
\hline
Seleksjonsbias & Risiko når man måler ``appbrukere handler oftere'' --- kanskje
fordi storkunder laster ned appen, ikke fordi appen får dem til å handle mer. \\
\hline
Brugsforeninger & Lokale, kooperative foreninger i bredere Coop-familie
utenfor Coop Danmark A/S. \\
\hline
EBITDA / EBIT & EBITDA = drift før av- og nedskrivninger. EBIT = drift etter.
Coop EBITDA: $+313$ mDKK; EBIT: $-215$ mDKK. Ikke bland! \\
\hline
\end{tabular}

\subsection*{G. Analytiske grep og fallgruver}

\begin{tabular}{p{4cm}p{12cm}}
\hline
\textbf{Begrep} & \textbf{Forklaring} \\
\hline
Value creation\newline $\neq$ value capture & Rød tråd: man kan skape
kundeverdi uten å fange økonomisk verdi. Oppgavens analytiske kjerne. \\
\hline
Outputs $\neq$ outcomes & Aktivitet er ikke det samme som resultat. Brukes mot
Lobyco-tallene. \\
\hline
Korrelasjon $\neq$\newline kausalitet & Lobyco-claim om ``handler oftere'' er
korrelasjon --- ikke bevis på at appen \emph{forårsaker} mer handel. \\
\hline
Strategisk spenning\newline (Coop $\leftrightarrow$ Lobyco) & Coop App =
omnichannel-instrument; Lobyco = potensiell modular producer. To
DBM-arketyper i samme organisasjon. \\
\hline
Mønsterslutn,\newline ikke kausalitet & Konkurrentsammenligningen viser at
digitale verktøy passer best når de matcher forretningsmodellen --- ikke
bevist årsakssammenheng. \\
\hline
\end{tabular}

\subsection*{H. Hurtigreferanse --- forfattere}

\begin{tabular}{p{4.5cm}p{11.5cm}}
\hline
\textbf{Forfatter(e)} & \textbf{Kjerne i én linje} \\
\hline
Duus \& Cooray (2020) & DVC --- relevance-sentrert digital kundeverdi. \\
\hline
Weill \& Woerner (2018) & DBM --- fire forretningsmodell-arketyper for digital
tid. \\
\hline
Wessel et al. (2021) & DT vs. ITE --- digital teknologi redefinerer eller
støtter identitet. \\
\hline
Porter (1996) & Strategi = posisjonering, ikke driftseffektivitet. \\
\hline
Osterwalder \& Pigneur (2010) & Business Model Canvas --- 9 byggesteiner. \\
\hline
Yoo et al. (2010) & Digital innovasjon = rekombinasjon av komponenter. \\
\hline
Nambisan et al. (2017) & Digital teknologi i selve innovasjonsprosessen. \\
\hline
Hukal \& Henfridsson (2017) & Samskapelse via rekombinasjon. \\
\hline
Rogers (1962/1983) & Diffusjon --- innovasjon oppleves som ny. \\
\hline
Christensen & Disruptiv vs. sustaining innovasjon. \\
\hline
Hedman \& Bjørn-Andersen (2016) & Coops IT-legacy og teknologisk gjeld. \\
\hline
Tawse \& Tabesh (2021) & 3C --- competency, commitment, coordination. \\
\hline
\end{tabular}
